\documentclass[11pt]{article}
\usepackage[margin=0.9in]{geometry}
\usepackage{amsmath,amssymb,amsthm,booktabs}
\usepackage[T1]{fontenc}
\usepackage{lmodern,microtype}
\usepackage[colorlinks=true,linkcolor=blue,urlcolor=blue]{hyperref}
\newtheorem{theorem}{Theorem}[section]
\newtheorem{lemma}[theorem]{Lemma}
\newtheorem{proposition}[theorem]{Proposition}
\newtheorem{corollary}[theorem]{Corollary}
\theoremstyle{remark}
\newtheorem{remark}[theorem]{Remark}
\newcommand{\norm}[1]{\left\|#1\right\|}
\title{Toward a four-dimensional Yang--Mills construction\\
\large Notebook 14: nonflat backgrounds and energy stability}
\author{Research working document}
\date{16 October 2026}
\setlength{\emergencystretch}{5em}
\begin{document}
\maketitle
\begin{abstract}
We pass from fixed-reference nonlinear averages to the actual
moving-reference averaging operators in a small global connection
chart. An explicit endpoint term corrects the change of derivative
coordinates. An exact cumulative gauge relation controls ordinary
reference backgrounds uniformly through the blocking depth, giving
small perturbations of both source averaging operators in physical
weighted norms. A separate Wilson plaquette calculation and a form
resolvent argument then control a gauge-fixed inverse, including the
scalar compression that defines its gauge projection. The inverse
estimates concern the energy blocks, with at most one derivative on
each side. They apply on a small global physical C1 ball, not yet on
the source's full local-frame and retained-history domains.
\end{abstract}
\noindent\textbf{Status.}
Research proofs checked through separate internal derivations and
audits. No external certification or novelty claim is made. The
explicit source-coordinate matches below refer to the deposited
transcriptions at the pinned commit, not to newly recovered original
CMP98 or CMP99b papers. The action and transpose normalizations are
stated explicitly at the point where the Wilson operator is used.
% Notebook 14 body fragment, no preamble. Unique label prefix: nc-.
% Requires amsmath, amssymb, amsthm and a one-argument \norm macro.

\section{The exact change of reference in the averaging differential}
\label{nc-section}

The derivative of the corrected average with fixed reference one
is not the same coordinate operator as the source average $Q(V)$
based at a nonflat reference $V$. We now compute their relation,
including the endpoint term. The block ratio $L$ is fixed, all
logarithms use their common identity branch, and all Hilbert norms
use physical volume. This calculation is internally audited at the
level of the pinned deposited transcriptions; no fresh verification
of the original CMP98 typesetting is claimed.

\subsection{Two references and an exact group identity}
\label{nc-group}

Write $M$ for the ordinary group average. Raw I:46965--46970
defines
\[
 \Phi^{\mathfrak g}(V,B;y,x)=\frac1i\log[
   (e^{iB}V)\langle y,x\rangle V\langle y,x\rangle^{-1}],
 \qquad
 \Phi^{\flat}(V,B;c)=\frac1i\log[M(e^{iB}V)(c)M(V)(c)^{-1}].
\]
For $U$ near $V$ set
\begin{equation}
 g_y^V(U)=\exp\left(i\sum_{x\in B(y)}L^{-d}
       \frac1i\log[U\langle y,x\rangle V\langle y,x\rangle^{-1}]
                   \right).
 \label{nc-endpoint-mean}
\end{equation}
The bridge I:47076--47087, expressly incorporating CMP98 (120),
identifies the corrected group increment as
\begin{equation}
 C_V(U)(c)=g_{c_-}^V(U)^{-1}M(U)(c)g_{c_+}^V(U)M(V)(c)^{-1}.
 \label{nc-moving-corrected}
\end{equation}
It equals one at $U=V$. With the logarithmic/linear split
I:24344--24349, the same bridge gives
\begin{equation}
 C_V(e^{iB}V)=e^{i\mathcal Q(V,B)},\qquad
 \mathcal Q(V,B)=LQ(V)B+O(B^2).
 \label{nc-source-linearization}
\end{equation}
This is the deposited identification of the middle coordinate;
the earlier transcription does not itself reproduce CMP98 (120).
We use the explicit bridge, rather than guessing that omitted input
formula.

Set $a_y(U)=g_y^1(U)$. The corrected map with fixed identity
reference is
\begin{equation}
 C_1(U)(c)=a_{c_-}(U)^{-1}M(U)(c)a_{c_+}(U).
 \label{nc-fixed-corrected}
\end{equation}
For fixed $V$ define
$k_y(U;V)=g_y^V(U)^{-1}a_y(U)a_y(V)^{-1}$. Direct substitution
and cancellation, preserving the factor order, give
\begin{equation}
 \begin{split}
 C_1(U)(c)C_1(V)(c)^{-1}
 =\operatorname{Ad}_{a_{c_-}(V)^{-1}}\bigl[
 k_{c_-}(U;V)^{-1}C_V(U)(c)
       \operatorname{Ad}_{M(V)(c)}k_{c_+}(U;V)\bigr].
 \end{split}
 \label{nc-group-identity}
\end{equation}
In particular no multiplicative cocycle for the endpoint mean has
been assumed.

For $U(s)=e^{isB}V$ let
\[
 \eta_y^1(V;B)=\frac1i\left.\frac d{ds}\right|_0
       a_y(U(s))a_y(V)^{-1},\qquad
 \eta_y^V(B)=\frac1i\left.\frac d{ds}\right|_0g_y^V(U(s)),
 \qquad Z_VB(y)=\eta_y^1(V;B)-\eta_y^V(B).
\]
Define the \emph{unscaled} endpoint difference by
\begin{equation}
 d_Wz(c)=z(c_-)-\operatorname{Ad}_{W(c)}z(c_+).
 \label{nc-endpoint-difference}
\end{equation}
This explicit definition fixes the sign independently of other
covariant-derivative conventions. All middle factors in
\eqref{nc-group-identity} equal one at $s=0$, so differentiating
gives
\begin{equation}
 \frac1i\left.\frac d{ds}\right|_0
 C_1(e^{isB}V)(c)C_1(V)(c)^{-1}
 =\operatorname{Ad}_{a_{c_-}(V)^{-1}}
       [LQ(V)B(c)-d_{M(V)}(Z_VB)(c)].
 \label{nc-angular-derivative}
\end{equation}

As a separate check, I:46977 defines
$LQ^{\flat}(V)B=\Phi_1^{\flat}(V;B)$. Differentiating the
general bridge directly at $B=0$ gives
\begin{equation}
 LQ(V)B=LQ^{\flat}(V)B-d_{M(V)}\overline r(V;B),
 \qquad
 \overline r(V;B)(y)=\sum_xL^{-d}\Phi_1^{\mathfrak g}(V;B;y,x).
 \label{nc-rooted-versus-corrected}
\end{equation}
Thus the later rooted operator $Q^{\flat}$ and the corrected source
operator $Q$ are distinct. Substitution into
\eqref{nc-angular-derivative} reproduces the direct derivative of
\eqref{nc-fixed-corrected}, confirming its endpoint sign.

\subsection{An explicit formula for the endpoint correction}
\label{nc-jacobians}

For a small complex Lie-algebra element $X$ put
\begin{equation}
 \mathcal J_XK=\int_0^1\operatorname{Ad}_{e^{itX}}K\,dt,
 \qquad
 \frac1i\left.\frac d{ds}\right|_0e^{i(X+sK)}e^{-iX}
           =\mathcal J_XK.
 \label{nc-exponential-jacobian}
\end{equation}
The integral power series and a Neumann series show that
$\mathcal J_X$ and its inverse are analytic near zero, and
$\norm{\mathcal J_X-I}+\norm{\mathcal J_X^{-1}-I}\le C\norm{X}$.
This does not require unitarity.

Write
\[
 V\langle y,x\rangle=e^{iY_x},\qquad
 \overline Y_y=\sum_xL^{-d}Y_x,\qquad
 r_x(V;B)=\Phi_1^{\mathfrak g}(V;B;y,x).
\]
The last quantity includes the actual contour orientations,
transports and repetitions. Applying
\eqref{nc-exponential-jacobian} to each contour logarithm and then
to the exponential of their mean proves
\begin{equation}
 Z_VB(y)=\mathcal J_{\overline Y_y}
       \sum_xL^{-d}\mathcal J_{Y_x}^{-1}r_x(V;B)
          -\sum_xL^{-d}r_x(V;B).
 \label{nc-Z-formula}
\end{equation}
In particular $Z_1=0$. The leading background term is
\[
 Z_VB=\frac i2\left(
   [\overline Y,\overline r]-\sum_xL^{-d}[Y_x,r_x]\right)
   +O\bigl((\max_x|Y_x|)^2\max_x|r_x|\bigr).
\]
There is no general identity forcing this difference of means of
commutators to vanish. The endpoint correction in
\eqref{nc-angular-derivative} must therefore be retained.

\subsection{The physical derivative conversion and its bound}
\label{nc-physical}

At fine spacing $h$ write $V=e^{ihA}$ and define the fixed-reference
physical map by
$C_1(e^{ihA})(c)=e^{iLh\mathcal F_h(A)(c)}$. For an ordinary
potential variation $K$, the fine angular left increment is
$h\mathcal J_{hA}K$ and the coarse left increment is
$Lh\mathcal J_{Lh\mathcal F_h(A)}D\mathcal F_h(A)K$.
Substituting these into \eqref{nc-angular-derivative} gives
\begin{equation}
 \begin{split}
 \mathcal J_{Lh\mathcal F_h(A)}D\mathcal F_h(A)K
 =\operatorname{Ad}_{a_-(V)^{-1}}\bigl[
 Q_h(V)\mathcal J_{hA}K
  -L^{-1}d_{M(V)}Z_V(\mathcal J_{hA}K)\bigr].
 \end{split}
 \label{nc-physical-unrearranged}
\end{equation}
Consequently the exact operator conversion is
\begin{equation}
 \boxed{\displaystyle
 Q_h(V)=\operatorname{Ad}_{a_-(V)}
   \mathcal J_{Lh\mathcal F_h(A)}D\mathcal F_h(A)
        \mathcal J_{hA}^{-1}+L^{-1}d_{M(V)}Z_V.}
 \label{nc-physical-conversion}
\end{equation}
Here $Q_h(V)$ is the normalized source coefficient placed between
the physical spaces at $h$ and $Lh$. Neither $d_{M(V)}$ nor $Z_V$
contains an inverse mesh factor. The argument of $Z_V$ before
division is $h\mathcal J_{hA}K$; its linearity cancels $h$ against
the coarse $Lh$ and leaves exactly $1/L$.

\begin{proposition}[Actual one-step source comparison]
\label{nc-one-step-proposition}
For fixed $L$, if $h\norm{A}_\infty$ lies in a sufficiently small
common complex neighborhood, then
\begin{equation}
 \norm{Q_h(e^{ihA})-Q_h(1)}_{p,h\to p,Lh}
     \le C_{L,p}h\norm{A}_\infty,
 \qquad p=2,\infty.
 \label{nc-one-step-bound}
\end{equation}
The constants are independent of $h$, refinement and compatible
periodic volume. The same estimate holds for absolute coefficient
row sums and physical-adjoint row sums, after changing the fixed
constant.
\end{proposition}

\begin{proof}
Each one-step contour has bounded length at fixed $L$. Products of
possibly nonunitary complex transports therefore give
$\max_x|Y_x|+|\overline Y|\le C_Lh\norm{A}_\infty$ and
$\norm{r_x(V;\cdot)}_{\infty\to\mathrm{color}}\le C_L$.
Equation \eqref{nc-Z-formula} implies
$\norm{Z_V}_{\infty\to\infty}\le C_Lh\norm{A}_\infty$.

For the physical Hilbert bound expand $Z_V$ into its finitely many
one-step bond coefficients. Each is analytic, vanishes at $V=1$,
and is bounded by $C_Lh\norm{A}_\infty$. Rows and columns have
fixed finite incidence; the physical adjoint supplies the fixed
volume ratio $L^d$. The Schur bound thus gives the same order
from fine physical $L^2$ to coarse-site physical $L^2$, with no
macroscopic path-trace estimate. The endpoint map $d_{M(V)}$ has
two entries per row, fixed column incidence and bounded complex
adjoints, so it is uniformly bounded in both norms.

The other coordinate factors in \eqref{nc-physical-conversion}
differ from the identity by at most $C_Lh\norm{A}_\infty$.
Combine this with the proved one-step estimate
$\norm{D\mathcal F_h(A)-Q_h(1)}_{p\to p}
\le D_ph\norm{A}_\infty$ and the flat contractions to obtain
\eqref{nc-one-step-bound}. Equivalently, the entire corrected
source coefficient in \eqref{nc-source-linearization} is a finite
local analytic function of the link logarithms $hA$, equal to
$Q_h(1)$ at zero. Its coefficientwise Taylor bound and the same
finite incidence prove the row and physical-adjoint row assertions
directly. Formula \eqref{nc-physical-conversion} records the
coordinate factors hidden by that shorter proof. Periodic
multiplicity counting preserves all constants.
\end{proof}

For a multilevel source comparison this estimate must be evaluated
at the actual ordinary reference backgrounds $M^t(U)$, whose
amplitudes are controlled in the following argument. It does not
license replacing them by the corrected fixed-reference iterates.
It also proves no nonflat inverse or minimizing-field theorem by
itself.

% Notebook 14 body fragment; no standalone preamble.
% Requires amsmath/amssymb/amsthm and a one-argument \norm macro.
% Unique label prefix: nb-. Independently audited research proof.

\section{Ordinary reference backgrounds and their averaging operators}
\label{nb-section}

The corrected iteration at reference $1$ and the ordinary group
average have different flat linearizations. We relate the complete
configurations by an explicit gauge, then estimate the actual
moving-reference source operators at the ordinary backgrounds.
This argument does not equate those operators with the ordinary
derivative of the fixed-reference corrected iteration.

Fix odd $L\ge3$, $M=L^j$, $\eta=M^{-1}$ and
\begin{equation}\label{nb-scales}
 h_t=L^t\eta\quad(0\le t\le j),\qquad
 \sum_{s<t}h_s=\frac{h_t-\eta}{L-1}.
\end{equation}
We use compatible periodic lattices or the full lattice, globally
defined complex Lie-algebra bond fields, and physical volume in
every Hilbert space. All constants below are independent of $j$,
$M$ and compatible periodic volume. They may depend on fixed
$L$, dimension, color norm and the averaging prescription.
Matrix estimates use a submultiplicative norm; fixed norm
equivalences can be absorbed in the constants.

\subsection{The exact endpoint identity}
\label{nb-source}

Let $\mathsf M$ denote the ordinary group-valued average. For a
configuration $V$ near the identity put
\begin{equation}\label{nb-endpoint-map}
 \alpha(V;y)=L^{-d}\sum_{x\in B(y)}\frac1i
                      \log V(\Gamma_{y,x}),\qquad
 a(V;y)=e^{i\alpha(V;y)}.
\end{equation}
The corrected group map at reference $1$ is
\begin{equation}\label{nb-corrected-map}
 \mathsf C(V)(c)
   =a(V;c_-)^{-1}\mathsf M(V)(c)a(V;c_+).
\end{equation}
This is the specialization of deposited Raw I:24313--24325,
equivalently its explicit coordinate bridge I:47082--47086.
For a general reference the endpoint factor also carries an
adjoint by the average of that reference. It disappears here
because the reference is $1$.

Use the gauge convention
$V^g(c)=g(c_-)V(c)g(c_+)^{-1}$. Ordinary gauge covariance is
\begin{equation}\label{nb-gauge-covariance}
 \mathsf M(V^g)=\mathsf M(V)^{g|_{\rm coarse}}.
\end{equation}
The deposit records restriction to base points at
I:43018--43026, and expressly invokes complexified gauge
covariance from CMP98 (11) at I:82901--82904. In the common
identity neighborhood it also follows from similarity covariance
of the analytic products, logarithms and means. All gauges used
below stay inside that neighborhood.

These are source matches to the pinned deposited transcription.
The full original CMP98 paper has not been independently
reverified in this calculation. In particular, the bridge's
explicit incorporation of its omitted coordinate (120) is kept
as the source convention, rather than replaced by a guessed
coordinate formula.

\subsection{Uniform ordinary connection amplitudes}

Let $U=e^{i\eta A}$ and $a=\norm{A}_\infty$. Write
\begin{equation}\label{nb-two-trajectories}
 C_0=U,\quad C_{t+1}=\mathsf C(C_t),\qquad
 V_t=\mathsf M^t(U).
\end{equation}
The fixed-reference theorem of Notebook 13 gives a radius
$\rho>0$, physical corrected fields $A_t$ and a constant
$D_\infty$, such that for $a<\rho$
\begin{equation}\label{nb-corrected-input}
 C_t=e^{ih_tA_t},\qquad \norm{A_t}_\infty\le2a,
 \qquad A_{t+1}=\mathcal F_{h_t}(A_t),
\end{equation}
and
\begin{equation}\label{nb-one-step-derivative}
 \norm{D\mathcal F_h(B)}_{\infty\to\infty}
             \le1+D_\infty h\norm{B}_\infty.
\end{equation}
The derivative bound is used on a common smaller ball, which
contains the corrected trajectories and their straight-line
interpolations with translates.

\begin{theorem}\label{nb-trajectory-theorem}
There are $a_*>0$ and $C_V<\infty$, independent of refinement
and volume, such that if $\norm{A}_\infty<a_*$, every ordinary
reference $V_t$ has a holomorphic physical logarithm
\begin{equation}\label{nb-ordinary-amplitude}
 B_t=(ih_t)^{-1}\log V_t,\qquad
             \norm{B_t}_\infty\le C_V\norm{A}_\infty.
\end{equation}
Moreover there is an explicit vertex gauge $g_t$ with
$V_t=C_t^{g_t}$ and $g_t^{\pm1}-1=O(h_t\norm{A}_\infty)$.
All these maps read a bounded final-scale neighborhood.
\end{theorem}

\begin{proof}
Let $m_\Gamma$ be the largest number of oriented bonds in a
one-step contour. It is fixed when $L$ is fixed; the conservative
bound $m_\Gamma\le dL$ suffices. Decrease $a_*$ so that
$2m_\Gamma a_*\le\log(3/2)$. By
\eqref{nb-corrected-input},
\[
 \norm{C_t(\Gamma)-1}
      \le e^{2m_\Gamma h_ta}-1\le\tfrac12.
\]
The logarithm series and the bound
$e^s-1\le3s/2$ for $0\le s\le\log(3/2)$ give, with slack,
\begin{equation}\label{nb-alpha-bound}
 \norm{\alpha(C_t;y)}\le K_\Gamma h_ta,
                   \qquad K_\Gamma=8m_\Gamma.
\end{equation}
Set
\begin{equation}\label{nb-gauge-recursion}
 g_0=1,\qquad
 g_{t+1}(y)=g_t(y)a(C_t;y),\qquad y\in\Lambda_{t+1}.
\end{equation}
The lattices are nested, so the restriction of $g_t$ is defined.
If $V_t=C_t^{g_t}$, then \eqref{nb-gauge-covariance} and
\eqref{nb-corrected-map} imply
\[
 V_{t+1}(c)=g_t(c_-)a(C_t;c_-)C_{t+1}(c)
                       a(C_t;c_+)^{-1}g_t(c_+)^{-1}.
\]
Thus $V_t=C_t^{g_t}$ by induction, with exactly the product
order in \eqref{nb-gauge-recursion}.

The ordered exponential products and
\eqref{nb-scales}--\eqref{nb-alpha-bound} give
\begin{align}
 \norm{g_t^{\pm1}}
       &\le\exp\left(\frac{K_\Gamma a(h_t-\eta)}{L-1}\right),
       \label{nb-gauge-size}\\
 \norm{g_t^{\pm1}-1}
       &\le\exp\left(\frac{K_\Gamma a(h_t-\eta)}{L-1}\right)-1.
       \label{nb-gauge-difference}
\end{align}
The second bound follows by telescoping; no unitarity or
commutativity is required. Put
\begin{equation}\label{nb-explicit-constants}
 K_V=2+\frac{2K_\Gamma}{L-1},\qquad C_V=3K_V.
\end{equation}
The same product estimate for
$V_t(c)=g_t(c_-)e^{ih_tA_t(c)}g_t(c_+)^{-1}$ yields
\begin{equation}\label{nb-bond-near-identity}
 \norm{V_t(c)-1}\le e^{K_Vah_t}-1.
\end{equation}
Take also $K_Va_*\le\log(3/2)$. The analytic logarithm then
satisfies
\[
 \norm{\log V_t(c)}
    \le2(e^{K_Vah_t}-1)\le3K_Vah_t,
\]
which proves \eqref{nb-ordinary-amplitude}.

Further reduce $a_*$, if necessary, to place the right side of
\eqref{nb-bond-near-identity} in the fixed one-step analytic
neighborhood of $\mathsf M$. This also justifies the construction
inductively before each ordinary average is applied. Holomorphy
follows from the fixed-reference theorem and the explicit
products and logarithms. Each one-step dependency has physical
radius at most a fixed multiple of $h_t$, so the accumulated
radius is bounded by a multiple of the geometric sum
\eqref{nb-scales}. Alternatively, use the deposited two-final-block
support of the corrected bond iterate.
\end{proof}

\subsection{Physical first differences and curvature}

Supremum smallness alone does not give bounded physical
curvature. In a commuting direction take $A_2$ alternating
between $a$ and $-a$ at neighboring sites in direction $1$,
with all other components zero. The plaquette angle is of
order $\eta a$, hence the curvature divided by $\eta^2$ is
of order $a/\eta$. This obstruction is already present at $t=0$.

\begin{theorem}\label{nb-C1-theorem}
In addition to Theorem~\ref{nb-trajectory-theorem}, assume
\begin{equation}\label{nb-input-gradient}
 b=\max_\nu\norm{\nabla_{\eta,\nu}A}_\infty<\infty.
\end{equation}
Then, with fixed constants,
\begin{align}
 \max_\nu\norm{\nabla_{h_t,\nu}A_t}_\infty
       &\le e^{2D_\infty a/(L-1)}b,
                                      \label{nb-corrected-C1}\\
 \max_\nu\norm{\nabla_{h_t,\nu}B_t}_\infty
       &\le C e^{Ca}b,                  \label{nb-ordinary-C1}\\
 h_t^{-2}\norm{V_t(\partial p)-1}
       &\le C e^{Ca}(b+a^2).            \label{nb-curvature}
\end{align}
These estimates hold for real or complex fields on compatible
periodic lattices, with the same constants.
\end{theorem}

\begin{proof}
The fixed-block map $\mathcal F_h$ commutes with translations
by $Lh$ in each coordinate direction. For such a translation
$T$, integrate its derivative along the segment from $B$ to
$TB$. This segment remains in the same supremum ball. Equation
\eqref{nb-one-step-derivative} gives
\[
 \norm{\mathcal F_h(TB)-\mathcal F_h(B)}_\infty
  \le(1+D_\infty h\norm{B}_\infty)\norm{TB-B}_\infty.
\]
Telescoping $L$ fine steps shows
$\norm{TB-B}_\infty/(Lh)\le
\norm{\nabla_{h,\nu}B}_\infty$. Iterating, using
\eqref{nb-corrected-input} and \eqref{nb-scales}, proves
\eqref{nb-corrected-C1}. This requires only coarse translation
covariance, not invariance under an arbitrary fine translation.

On the same fixed coordinate ball, differentiation of the
finite contour formula gives
\begin{equation}\label{nb-contour-derivative}
 \norm{D_A\alpha(e^{ihA})}_{\infty\to\infty}
 +\norm{D_Aa(e^{ihA})}_{\infty\to\infty}\le C_\Gamma h.
\end{equation}
Indeed, each differentiated bond exponential contributes $h$;
the other factors have bounded product norm. The logarithm
derivative is bounded on $\norm{V-1}\le1/2$, as is the final
exponential derivative. No fine derivative is taken here.

For $s<t$, translation by $h_t$ preserves the scale-$s$ block
construction. By \eqref{nb-contour-derivative}, telescoping
$h_t/h_s$ input steps and using \eqref{nb-corrected-C1},
\begin{equation}\label{nb-endpoint-spatial-difference}
 \norm{a(C_s;y+h_te_\nu)-a(C_s;y)}
        \le C h_sh_t e^{2D_\infty a/(L-1)}b.
\end{equation}
The inverse factors obey the same estimate. Telescope the
ordered products defining $g_t$ at the two vertices, bound the
undifferentiated products by \eqref{nb-gauge-size}, and sum
$h_s$. This gives
\begin{equation}\label{nb-gauge-spatial-difference}
 \norm{g_t(y+h_te_\nu)^{\pm1}-g_t(y)^{\pm1}}
                         \le C e^{Ca}h_t^2b.
\end{equation}
The corresponding difference of $C_t(c)=e^{ih_tA_t(c)}$
is bounded by the same expression. The identity $V_t=C_t^{g_t}$
therefore bounds the difference of $V_t$ across $h_te_\nu$
by $Ce^{Ca}h_t^2b$. The logarithm is uniformly Lipschitz in
the half-unit ball, by integrating its derivative on the convex
segment. Dividing the logarithm difference by $h_t^2$ proves
\eqref{nb-ordinary-C1}. The case $t=0$ is the input estimate.

For curvature, first work with $C_t$. Pairing the two parallel
sides of a plaquette, the difference formula for exponentials
bounds the opposite-side differences by
$Ch_t^2\max_\nu\norm{\nabla_{h_t,\nu}A_t}_\infty$.
The remaining constant-side product is a group commutator,
bounded by $Ch_t^2a^2$. Thus
\[
 \norm{C_t(\partial p)-1}
 \le Ch_t^2\left(e^{2D_\infty a/(L-1)}b+a^2\right).
\]
The ordinary plaquette is conjugated by $g_t$ at its initial
vertex. Equation \eqref{nb-gauge-size} proves
\eqref{nb-curvature}. For real fields the gauges are unitary,
so a unitarily invariant norm avoids this conjugation loss.
Periodization preserves the translation identities and all
telescoping estimates, including on short compatible tori.
\end{proof}

\subsection{Actual bond averages in physical norms}

Let $Q_h(V)$ be the normalized linear coefficient of the
corrected relative increment at reference $V$, held fixed.
It is the one-step source operator of CMP98 (124), transcribed
at Raw I:24355--24369. The explicit bridge I:47082--47086
expresses it through finite local analytic operations. At $V=1$
it is the straight-string operator $Q_h^{\rm st}$.

For an operator from spacing $h$ to spacing $H$, define the
absolute physical Schur norm by
\begin{equation}\label{nb-av-norm}
 \norm{P}_{\rm av}
 =\max\left\{
       \sup_c\sum_b\norm{P_{cb}},\quad
       \sup_b(H/h)^d\sum_c\norm{P_{cb}}\right\}.
\end{equation}
This controls the bounded-field norm, the bounded-field norm
of the physical adjoint and the physical $L^2$ norm. It is
submultiplicative across scales: the volume ratios multiply.
For the flat normalized straight-string map both row sums
equal one. In fact each input bond occurs $L$ times among
the $L^d$ translated strings of length $L$, so its ordinary
column sum is $L^{-d}$ and its physical adjoint column sum
is one.

Finite-dimensional local analyticity and fixed-$L$ row and
column incidence give constants $r_Q,C_Q>0$ with
\begin{equation}\label{nb-local-Q-difference}
 \norm{Q_h(e^{iX})-Q_h^{\rm st}}_{\rm av}
             \le C_Q\norm{X}_\infty,
                 \qquad\norm{X}_\infty<r_Q.
\end{equation}
To verify this directly, differentiate each local coefficient
on a smaller closed complex polydisc and sum its finitely many
incidences. The physical column factor is $L^d$, independent
of $h$. The separate exact change-of-reference calculation
identifies the logarithm Jacobians and endpoint correction
that enter these coefficients; this shorter norm argument
uses their already explicit local analytic definition.

The actual source product uses the ordinary backgrounds:
\begin{equation}\label{nb-actual-product}
 Q_j(U)=Q_{h_{j-1}}(V_{j-1})\cdots Q_\eta(V_0).
\end{equation}
This is the composition recorded at Raw I:25147--25155 and
I:47046--47061. By Theorem~\ref{nb-trajectory-theorem},
$\norm{\log V_t}_\infty\le C_Vah_t$. Reduce $a_*$ so that
$C_Va_*<r_Q$. Submultiplicativity and an exact telescoping
of the difference of products prove
\begin{align}
 \norm{Q_j(U)}_{\rm av}
       &\le\exp\left(\frac{C_QC_Va}{L-1}\right),
                                      \label{nb-product-bound}\\
 \norm{Q_j(U)-Q_M^{\rm st}}_{\rm av}
       &\le\exp\left(\frac{C_QC_Va}{L-1}\right)-1.
                                      \label{nb-product-difference}
\end{align}
For a partial product ending at $h_t$, replace $(L-1)^{-1}$
in the exponent by $(h_t-\eta)/(L-1)$. These are actual source
operator estimates, obtained without renaming
$D\mathcal F^{(j)}(A)$ as $Q_j(U)$.

\subsection{Scalar averages and finite physical range}

The scalar average called $Q'_j$ in the deposit, and $Q_0$
in the homogeneous flat analysis, is
\begin{equation}\label{nb-scalar-average}
 (Q'_j(U)\lambda)(y)=L^{-jd}\sum_{x\in B^j(y)}
       \operatorname{Ad}_{U(\Gamma^{(j)}_{y,x})}\lambda(x).
\end{equation}
Raw I:23913--23920 states this formula. Its contour follows
successive block centers. The description at I:17798--17810
uses level-$s$ ordinary averaged configurations on the legs,
with at most $dL$ bonds per leg. The ordinary amplitude bound
and \eqref{nb-scales} therefore imply
\begin{equation}\label{nb-scalar-transport}
 \norm{U(\Gamma^{(j)}_{y,x})^{\pm1}}\le e^{Cah_j},
 \qquad
 \norm{\operatorname{Ad}_{U(\Gamma^{(j)}_{y,x})}-I}
                                      \le e^{Cah_j}-1.
\end{equation}
At a fixed level the blocks partition the fine vertices.
Each input scalar vertex occurs in one row with weight
$L^{-jd}$, canceled by the physical adjoint volume factor.
Consequently
\begin{equation}\label{nb-scalar-difference}
 \norm{Q'_j(U)}_{\rm av}\le e^{Cah_j},\qquad
 \norm{Q'_j(U)-Q'_j(1)}_{\rm av}\le e^{Cah_j}-1.
\end{equation}
This scalar operator is distinct from the bond average and
from the symbol $Q_0$ used for the main term of CMP98 (125).

Partition the spaces into compatible physical cells of side
one, with their physical Hilbert measures. The finite
multilevel reading radius is bounded independently of
refinement. For bond averages the deposit states the sharper
two-final-block support at I:23927--23930 and I:24332--24336.
If $R$ bounds the physical reading radius and $N_R$ the
number of possible neighboring cells, every cell block has
operator norm at most the global Hilbert norm. Hence
\begin{equation}\label{nb-cell-Schur}
 \norm{P}_{\kappa,\mathrm{cell}}
 :=\max\left\{
 \sup_x\sum_y e^{\kappa d(x,y)}\norm{P_{xy}}_{2\to2},
 \sup_y\sum_x e^{\kappa d(x,y)}\norm{P_{xy}}_{2\to2}\right\}
 \le N_Re^{\kappa R}\norm{P}_{2\to2}
\end{equation}
for such a finite-range $P$. Applying this to the differences
in \eqref{nb-product-difference} and
\eqref{nb-scalar-difference} gives uniform $O(a)$ bounds in
the physical-cell weighted Schur norm. Periodic distance and
fixed cover multiplicities give the same result on small
compatible tori.

\subsection{Local holomorphy and scope}

All constructions above are holomorphic on the same complex
ball $\norm{A}_\infty<a_*$. At finite regulators this follows
from products of finite-dimensional holomorphic maps. Their
uniform bounds and bounded physical reading ranges also give
the corresponding bounded-field local maps on the full
lattice. In particular $A\mapsto Q_j(e^{i\eta A})$ and
$A\mapsto Q'_j(e^{i\eta A})$ are holomorphic in the weighted
operator norm just defined, with a uniform bound on every
smaller ball. For example, if that bound is $C_0$ and an
argument lies at distance at least $r$ from the boundary of
the chosen ball, the multivariable Cauchy formula gives
\begin{equation}\label{nb-Cauchy}
 \norm{D^nQ_j(A)[K_1,\ldots,K_n]}_{\kappa,\mathrm{cell}}
 \le C_0 n!\left(\frac er\right)^n
                         \prod_{m=1}^n\norm{K_m}_\infty,
\end{equation}
and the same statement for $Q'_j$. One obtains it using
polydisc radii $r/(n\norm{K_m}_\infty)$ and
$n^n\le e^n n!$. Directions of zero norm are handled by
multilinearity. This is a bounded-background derivative
statement; no Hilbert norm is substituted for a background
direction without another estimate.

The result supplies uniform ordinary reference trajectories,
their physical $C^1$ and curvature bounds when the initial
first differences are bounded, and actual source averaging
operators close to their flat counterparts. It concerns a
globally trivialized homogeneous class. The next section uses these estimates to control an explicitly
normalized Wilson Hessian and its moving projected gauge inverse.
The minimizing field is constructed separately in Notebook 15
on a small global physical $C^{1,\alpha}$ ball; Notebook 16
controls its nonlinear local responses and classical coarse action.
Mixed retained-history domains, strict background-support expansions
and an interacting Yang--Mills ultraviolet limit remain separate.
The original-primary-paper
verification limitation stated in Section~\ref{nb-source}
remains in force.

% Notebook 14 body fragment; no preamble. Unique prefix nw-.
% Requires amsmath, amssymb, amsthm and \norm; follows nc-/nb- fragments.
% Main proof independently audited; source action convention scope explicit.

\section{A nonflat Wilson Hessian and its full energy inverse}
\label{nw-section}

The averaging bounds above close the remaining averaging hypotheses
of an energy-form perturbation argument. We first calculate the
Wilson Hessian, including its normalization. Its bare action,
physical pairing and left-exponential chart match the deposited
conventions in dimension four. The resulting theorem uses the
deposited homogeneous averaging maps; the graded-curl and effective
Hessian identifications remain separate, as stated below.

\subsection{Action, charts, and an exact local decomposition}
\label{nw-wilson}

Write $\eta=M^{-1}\le1$ and $\tau=N_c^{-1}\operatorname{Tr}$.
Fine fields have physical pairing
$\langle B,C\rangle_\eta=\eta^d\sum_e\tau(B_eC_e)$.
The sum is over positive bonds; plaquettes below are counted once,
with $\mu<\nu$. On Hermitian fields this is the real Hilbert pairing;
its complex extension is bilinear. Superscript $T$ always denotes
that transpose, with no conjugation of a complex background.
Define
\begin{equation}\label{nw-action}
 \mathcal S_\eta(U)=\eta^{d-4}\sum_p
   \left(1-\frac12\tau(P_p+P_p^{-1})\right),\qquad
 U_e(s,t)=e^{i\eta(sB_e+tC_e)}U_e,
\end{equation}
and define $\Delta_U$ by the mixed derivative of this action.
There is no coupling prefactor in \eqref{nw-action}. At $U=1$,
$\Delta_1=d^Td$, where $d$ is the ordinary physical plaquette curl.

Let $U_e=e^{i\eta a_e}$, with
\begin{equation}\label{nw-sizes}
 \epsilon_0=\norm{a}_\infty,\qquad
 \epsilon_1=\norm{\nabla a}_\infty,\qquad
 F=\epsilon_1+\epsilon_0^2.
\end{equation}
The field can be complex; $\epsilon_0$ is bounded by a fixed small
constant. Comparing opposite plaquette sides by the integral formula
for matrix exponentials, and bounding the residual group commutator,
gives
\begin{equation}\label{nw-plaquette}
 \norm{P_p-I}+\norm{P_p^{-1}-I}\le C\eta^2F.
\end{equation}
Indeed opposite potentials differ by at most $\eta\epsilon_1$,
so their exponentials differ by $C\eta^2\epsilon_1$.
For the remaining two constant potentials the commutator costs
$C\eta^2\epsilon_0^2$. All path transports involved here have a
fixed number of links and norm at most $e^{C\eta\epsilon_0}$.

For the loop $e_1,\ldots,e_4$ around $p$, set
$V_i=U_{e_1}\cdots U_{e_{i-1}}$.
A negative traversal of the positive bond $f$ has exact chart
increment $B_{-f}=-\operatorname{Ad}_{U_f^{-1}}B_f$.
Put $X_i^B=\operatorname{Ad}_{V_i}B_{e_i}$ and
\begin{equation}\label{nw-curl}
 \mathcal C_U B=\eta^{-1}\sum_iX_i^B.
\end{equation}
The exact plaquette variation is
$P_p(s,t)=\prod_i e^{i\eta(sX_i^B+tX_i^C)}P_p$.
Every transport difference from identity is $O(\eta\epsilon_0)$,
and therefore
\begin{equation}\label{nw-curl-difference}
 \mathcal C_U=d+E_U,\qquad
 \norm{E_U}_{\infty\to\infty}
       +\norm{E_U^T}_{\infty\to\infty}\le C\epsilon_0.
\end{equation}
The transpose norm uses physical measure. Fixed plaquette incidence
also gives the same estimate in a weighted physical-cell Schur norm.

Writing $S_B=\sum_iX_i^B$, direct mixed differentiation gives
\begin{equation}\label{nw-mixed}
 \partial_s\partial_tP_p\big|_0=-\eta^2
 \left\{\frac{S_BS_C+S_CS_B}{2}
 +\frac12\sum_{i<j}
 ([X_i^B,X_j^C]+[X_i^C,X_j^B])\right\}P_p.
\end{equation}
For the oriented density $-\tau(P_p)$ the principal contribution is
$\eta^d\tau(((\mathcal C_UB)(\mathcal C_UC)
+(\mathcal C_UC)(\mathcal C_UB))P_p)/2$.
The remaining contribution is exactly
\begin{equation}\label{nw-commutator-remainder}
 \frac{\eta^d}{2\eta^2}\sum_{i<j}
 \tau\left(([X_i^B,X_j^C]+[X_i^C,X_j^B])(P_p-I)\right),
\end{equation}
because a commutator has trace zero. The reverse loop has curl
$-\operatorname{Ad}_{P_p^{-1}}\mathcal C_U$; cyclicity removes this
adjoint transport from its principal quadratic term. Thus
\begin{equation}\label{nw-hessian}
 \Delta_U=\mathcal C_U^TM_U\mathcal C_U+V_U,\qquad
 \tau(XM_{U,p}Y)=\frac14\tau((XY+YX)(P_p+P_p^{-1})),
\end{equation}
where $V_U$ is half \eqref{nw-commutator-remainder} plus half its
reverse-loop version. Projection onto the fixed traceless Lie space
is implicit in the definition of $M_U$. These coefficients are
bilinearly symmetric and holomorphic. Equation \eqref{nw-plaquette}
implies
\begin{equation}\label{nw-coefficients}
 \norm{M_U-I}\le C\eta^2F,
 \qquad \norm{V_U}_{\infty\to\infty}
       +\norm{V_U^T}_{\infty\to\infty}\le CF.
\end{equation}
Each curl has row and transpose-row norm $O(\eta^{-1})$.
Consequently the ordinary bounded local operator
$W_U=\Delta_U-\mathcal C_U^T\mathcal C_U$ has row and
transpose-row norm $CF$: the two inverse mesh factors are paid for
by $\eta^2$ in \eqref{nw-coefficients}. In particular
\begin{equation}\label{nw-small-form}
 \Delta_U-d^Td=d^TE_U+E_U^Td+E_U^TE_U+W_U
       =S_v^TV_{\rm kin}S_v,
 \quad S_v=(I,\nabla),
\end{equation}
with $\norm{V_{\rm kin}}_\kappa\le
C(\epsilon_0+\epsilon_1+\epsilon_0^2)$.
Here $\nabla$ contains all physical fine differences of bond
components; the curl is a fixed component combination of them.
The scalar gauge derivative is separately
\begin{equation}\label{nw-scalar-gradient}
 D_U\lambda(e)=\eta^{-1}
  (\operatorname{Ad}_{U_e}\lambda(e_+)-\lambda(e_-))
  =D\lambda+E_s\lambda,
 \qquad \norm{E_s}_\kappa\le C\epsilon_0.
\end{equation}
All coefficients in these two estimates have fixed fine range.

\subsection{The scalar compression and a form-resolvent lemma}
\label{nw-scalar}

For physical-cell block operators use the Banach algebra norm
\begin{equation}\label{nw-schur}
 \norm{T}_\kappa=\max\left\{
 \sup_x\sum_y e^{\kappa d(x,y)}\norm{T_{xy}},
 \sup_y\sum_x e^{\kappa d(x,y)}\norm{T_{xy}}\right\}.
\end{equation}
The norm is submultiplicative and invariant under transpose.
Fine-cell spaces have measure $\eta^d$, so it introduces no power
of the number of fine sites. The earlier flat inverse theorem gives
\begin{equation}\label{nw-flat-blocks}
 \norm{S_vG_0S_v^T}_\kappa\le K_v,
 \quad G_0=(d^Td+DR_0D^T+Q^Ta_{\rm pen}Q)^{-1},
\end{equation}
for $a_{\rm pen}$ in a fixed compact positive interval.
Both flat averages $Q=Q_M^{\rm st}$ and $Q_0$, and their transposes,
have uniform weighted norms by physical row/column contraction and
fixed final-scale range.

For clarity the additional scalar input is derived explicitly.
Let $G'_0=(\mathcal L+Q_0^Ta'Q_0)^{-1}$, $\mathcal L=D^TD$,
where $a'$ lies in a fixed compact positive interval, and let
$S_s=(I,D)$. In coarse Bloch aliases write $P_r$ for the fine
Laplacian multiplier, $a_r$ for the scalar average, and
\begin{align}
 A&=a_0a_0^-,&
 R_{s,k}&=\sum_{r\ne0}\frac{a_ra_r^-}{P_r^k},\nonumber\\
 D_{s,2}&=A+P_0^2R_{s,2},&
 N&=P_0+a'A+a'P_0R_{s,1}.                   \label{nw-scalar-denoms}
\end{align}
The scalar Woodbury formula has central entry
$(R_{s,1}+(a')^{-1})/[A+P_0(R_{s,1}+(a')^{-1})]$,
mixed entry $-a_0^-a_s/[P_s(A+P_0(R_{s,1}+(a')^{-1}))]$,
and a high diagonal $P_r^{-1}$ plus a rank-one correction.
These canceled expressions, high-alias gaps and alias Parseval
give $\norm{S_sG'_0S_s^T}_\kappa\le K_s$.
They also give the exact compression formula
\begin{equation}\label{nw-compression}
 C_0=Q_0(G'_0)^2Q_0^T=\frac{D_{s,2}}{N^2},
 \qquad C_0^{-1}=\frac{N^2}{D_{s,2}}.
\end{equation}
On an enlarged real Bloch cube, $A\ge c>0$ and the high remainders
are nonnegative and uniformly bounded. Thus $N,D_{s,2}$ are
bounded above and away from zero. Their uniformly bounded
holomorphic preliminary tube can be shrunk by a fixed amount to
retain these lower bounds. Periodic physical-cell reconstruction
and contour shifting prove
$\norm{C_0}_\kappa+\norm{C_0^{-1}}_\kappa\le K_c$.
The zero fiber gives $C_0=(a')^{-2}$; for $M=1$ it gives
$(P_0+a')^{-2}$. No inverse of the unpenalized zero mode is used.

We will repeatedly use the following elementary form identity.
If $G=\mathcal A^{-1}$, $K=SGS^T$ has weighted norm at most $K_*$,
and $I$ is a component of $S$, then for $K_*\norm{V}_\kappa<1$,
\begin{equation}\label{nw-form-resolvent}
 (\mathcal A+S^TVS)^{-1}
 =G-GS^TV(I+KV)^{-1}SG,
\end{equation}
and the augmented block has norm at most
$K_* /(1-K_*\norm{V}_\kappa)$ and differs from $K$ by at most
$K_*^2\norm{V}_\kappa/(1-K_*\norm{V}_\kappa)$.
The Neumann series in \eqref{nw-schur} proves these assertions and
holomorphy. Multiplication verifies the inverse at each fixed
spacing on the finite or infinite lattice. If the real baseline
$\mathcal A$ is positive definite, the same smallness gives
relative-form coercivity for real symmetric perturbations.

\subsection{Conditional closure of the projected gauge inverse}
\label{nw-conditional}

\begin{theorem}\label{nw-energy-theorem}
Suppose, in one common weighted algebra, that
\begin{gather}
 D_U=D+E_s,\quad Q_s(U)=Q_0+F_s,\quad Q_v(U)=Q+F_v,
 \quad \Delta_U=d^Td+S_v^TV_{\rm kin}S_v,\nonumber\\
 \norm{E_s}_\kappa+\norm{F_s}_\kappa+
 \norm{F_v}_\kappa+\norm{V_{\rm kin}}_\kappa\le\varepsilon.
                                                        \label{nw-hypotheses}
\end{gather}
Define
\begin{align}
 G'_U&=(D_U^TD_U+Q_s(U)^Ta'Q_s(U))^{-1},\nonumber\\
 C_U&=Q_s(U)(G'_U)^2Q_s(U)^T,\nonumber\\
 R_U&=I-G'_UQ_s(U)^TC_U^{-1}Q_s(U)G'_U,\nonumber\\
 \mathcal A_U&=\Delta_U+D_UR_UD_U^T
                     +Q_v(U)^Ta_{\rm pen}Q_v(U).       \label{nw-full-operator}
\end{align}
For a sufficiently small uniform $\varepsilon_*>0$, if
$\varepsilon\le\varepsilon_*$ all these inverses exist and
\begin{equation}\label{nw-energy-bound}
 \norm{S_sG'_US_s^T}_\kappa+
 \norm{C_U^{-1}}_\kappa+\norm{R_U}_\kappa+
 \norm{S_v\mathcal A_U^{-1}S_v^T}_\kappa\le C.
\end{equation}
Each displayed block differs from its flat counterpart by at
most $C\varepsilon$. Constants are uniform in refinement and
compatible periodic volume; the full-lattice statement also holds.
Holomorphic coefficient families give holomorphic inverse families.
On real data $R_U$ is the orthogonal projection off
$\operatorname{ran}(G'_UQ_s(U)^T)$.
\end{theorem}

\begin{proof}
The scalar precision difference is
\[
 D^TE_s+E_s^TD+E_s^TE_s+
 Q_0^Ta'F_s+F_s^Ta'Q_0+F_s^Ta'F_s.
\]
It equals $S_s^TV_sS_s$ with $\norm{V_s}_\kappa\le C_s\varepsilon$.
The form identity controls the scalar energy block and its
$O(\varepsilon)$ difference when $K_sC_s\varepsilon\le1/2$.
Product expansion gives $\norm{C_U-C_0}_\kappa\le C'_s\varepsilon$.
Equation \eqref{nw-compression} and
\[
 C_U^{-1}=(I+C_0^{-1}(C_U-C_0))^{-1}C_0^{-1}
\]
then control its inverse if $K_cC'_s\varepsilon\le1/2$.
The formula for $R_U$ gives its bound and its
$O(\varepsilon)$ difference. Its projection assertion is the Gram
projection formula; for complex coefficients it is a bilinear
holomorphic continuation rather than a positivity claim.

The gauge difference is exactly
\begin{equation}\label{nw-gauge-difference}
 D_UR_UD_U^T-DR_0D^T
 =D(R_U-R_0)D^T+E_sR_UD^T+DR_UE_s^T+E_sR_UE_s^T.
\end{equation}
The divergence $D^T$ on bonds is a fixed shifted component
combination of $\nabla$. Keeping these derivatives in $S_v$,
\eqref{nw-gauge-difference} is $S_v^TV_gS_v$ with norm
$C_g\varepsilon$. The average-penalty difference
$Q^Ta_{\rm pen}F_v+F_v^Ta_{\rm pen}Q+F_v^Ta_{\rm pen}F_v$
is a bounded $I,I$ coefficient of size $C_a\varepsilon$.
Thus $\mathcal A_U=\mathcal A_0+S_v^TV_{\rm tot}S_v$,
with $\norm{V_{\rm tot}}_\kappa\le C_v\varepsilon$.
A second application of \eqref{nw-form-resolvent}, with
$K_vC_v\varepsilon\le1/2$, proves the result.
\end{proof}

The theorem controls $G_U$, $\nabla G_U$, $G_U\nabla^T$ and
$\nabla G_U\nabla^T$, with exponential coarse-cell locality.
It does not propagate one-sided second derivatives by this
argument: differentiating $G_0S_v^T$ twice would require an
unavailable total third-order baseline block.

\subsection{Penalty independence and the homogeneous Wilson conclusion}
\label{nw-conclusion}

The scalar penalty is unnecessary for identifying the projection.
Fix $\mathcal L_U=D_U^TD_U$ and $Q_s$, and let
$G_b=(\mathcal L_U+Q_s^TbQ_s)^{-1}$,
$E_b=G_bQ_s^T$. Whenever the inverse and Gram inverse exist,
\begin{equation}\label{nw-range}
 \operatorname{ran}E_b
   =\{v:\mathcal L_Uv\in\operatorname{ran}Q_s^T\}.
\end{equation}
One inclusion follows by applying $\mathcal L_U$ to $G_bQ_s^T$.
For the other, $\mathcal L_Uv=Q_s^Tw$ implies
$v=G_bQ_s^T(w+bQ_sv)$.
The range is therefore independent of the symmetric coarse penalty
$b$. Its Gram inverse makes it nondegenerate, and
$E_b(E_b^TE_b)^{-1}E_b^T$ is the unique bilinear projection
onto this range along its orthogonal complement. Thus $R_U$ is
penalty independent. The identity holds directly for complex data
whenever the stated inverses exist, and on real data it is the
ordinary orthogonal projection identity.

We may consequently construct $R_U$ using $a'=1$.
It agrees with a differently penalized source formula wherever
that formula is defined, without estimating an unprinted scalar
penalty sequence. An estimate for that other scalar inverse
$G'_U$ itself still requires its own normalization bounds.

\begin{corollary}\label{nw-wilson-corollary}
Fix odd $L\ge3$, $M=L^j$, and the homogeneous full or periodic
lattice. Let $Q_v(U)$ be the deposited product of corrected
one-step derivatives at the ordinary reference trajectory, as
in \eqref{nb-actual-product}, and let $Q_s(U)$ be the deposited
scalar transported average \eqref{nb-scalar-average}.
Define $\Delta_U$ by \eqref{nw-action}, $D_U$ by
\eqref{nw-scalar-gradient}, and $R_U$ by
\eqref{nw-full-operator} with auxiliary scalar penalty one.
For $a_{\rm pen}$ in a fixed compact positive interval, there are
$\rho,\kappa,C>0$, independent of $j$ and compatible volume, such that
\begin{equation}\label{nw-final}
 \norm{a}_\infty+\norm{\nabla a}_\infty<\rho
 \quad\Longrightarrow\quad
 \norm{(I,\nabla)
 (\Delta_U+D_UR_UD_U^T+Q_v(U)^Ta_{\rm pen}Q_v(U))^{-1}
 (I,\nabla)^T}_\kappa\le C.
\end{equation}
The inverse is holomorphic on this complex global $C^1$ ball,
and its augmented block differs from the flat block by at most
$C(\norm{a}_\infty+\norm{\nabla a}_\infty)$.
\end{corollary}

\begin{proof}
Equations \eqref{nb-product-difference},
\eqref{nb-scalar-difference}, and \eqref{nb-cell-Schur}
give $\norm{Q_v(U)-Q}_\kappa+\norm{Q_s(U)-Q_0}_\kappa
\le C\epsilon_0$.
Their mechanism is the ordinary-reference estimate
$\norm{\log V_t}\le Ch_t\epsilon_0$, a physical one-step
coefficient difference $O(h_t\epsilon_0)$, and the geometric
sum $\sum_{t<j}h_t\le1/(L-1)$; no derivative of a different
nonlinear composition is substituted.
Equations \eqref{nw-small-form} and \eqref{nw-scalar-gradient}
supply the remaining hypotheses of Theorem~\ref{nw-energy-theorem}.
\end{proof}

The bare source normalization is matched at the deposited-transcription level.
Raw I lines 453--467 normalize $\operatorname{tr}1=1$, count
unoriented plaquettes once and define
$A(U)=\sum_p(1-\operatorname{Re}\operatorname{tr}P_p)$;
lines 472--481 place the coupling separately. Thus this bare action
is exactly \eqref{nw-action} in dimension four.
Lines 12756--12760 give the physical trace pairings with fine measure
$\eta^d$ and their bilinear transposes. Lines 15471--15478 and
15643--15648 define $\Delta$ as the quadratic coefficient in the
simultaneous left chart $\mathsf A(e^{i\eta A}U_0)$, exactly the chart
in \eqref{nw-action}. Lines 43103--43107 again specify the bare
Wilson action for constrained minimization. These passages resolve
the action-prefactor, physical-sum and chart conventions used here.
Lines 13583--13593 define the full source operator
$\Delta+D_URD_U^T+Q^TaQ$, with $a=1$ at unit scale.
At a nonzero global potential $A_0$, the ordinary-coordinate Hessian
of $A\mapsto\mathcal S_\eta(e^{i\eta A})$ also contains the background
current paired with the second derivative of the change to this left
chart. Matching the action does not remove that term or identify the
two Hessian matrices. The separate source claim about
$\Delta-D_U^TD_U$ still requires matching its precise graded
one-form covariant curl, and no mixed retained-history effective
Hessian is identified with this bare homogeneous Hessian.
These checks concern the pinned deposit, not a fresh verification
of the original journal proofs.

Corollary~\ref{nw-wilson-corollary} is a theorem for its explicitly
defined Wilson operator and deposited homogeneous averages.
It does not establish nonflat pointwise weak norms, nonflat current
domains, mixed retained histories, arbitrary nontrivial reference
connections, or an interacting continuum construction.

\begin{thebibliography}{9}
\bibitem{repo}
\href{https://github.com/ethanwsussman/Yang--Mills-LLM-experiment/tree/e392ef45013a581664bab76b90867170b30fe10e}
{Pinned Yang--Mills-LLM-experiment account}, commit
\texttt{e392ef4501\allowbreak 3a581664ba\allowbreak b76b908671\allowbreak 70b30fe10e}.
The exact locations and remaining source-identification checks are stated
alongside the corresponding arguments above.
\bibitem{previous}
Companion research proofs: Notebook 09, the full flat fine inverse
and scalar gauge identities; Notebook 12, pointwise and Schauder
regularity on the flat current domain; Notebook 13, the complete
fixed-reference nonlinear averaging ball and constraint chart.
\end{thebibliography}
\end{document}
